\documentclass[11pt,a4paper]{article}

\usepackage[margin=2.5cm]{geometry}
\usepackage{array}
\usepackage{booktabs}
\usepackage{etoolbox}
\usepackage{float}
\usepackage{graphicx}
\usepackage[hidelinks]{hyperref}
\AtBeginEnvironment{thebibliography}{\setlength{\itemsep}{0.65\baselineskip}}

\title{Topic 1: Issue Report Classification}
\author{Pham Ngoc Hai}
\date{}

\begin{document}

\maketitle

\begin{abstract}
The NLBSE'24 dataset contains 3,000 GitHub issue reports from five open-source
repositories. Each issue is labelled as a bug, feature, or question. This
project builds one classifier for each repository to predict these labels from
the issue title and description. We compare five text models and one dummy
baseline using five-fold cross-validation on the training set. Based on this
comparison, we choose word-and-character TF--IDF with LinearSVC. On the official
test set, our chosen model achieved a cross-repository macro-F1 of 0.754. We
also examined three exact train/test overlaps; excluding them changed macro-F1
by only 0.000257.
\end{abstract}

\section{Problem and Research Questions}

The NLBSE'24 tool competition ask us to classify GitHub issues as bugs, features, or questions. % INTENTIONAL ESL 01
Each model receive an issue title and description as input. % INTENTIONAL ESL 02
It returns exactly one label. The project answer three research questions. % INTENTIONAL ESL 03
These questions guides the model comparison and final evaluation. % INTENTIONAL ESL 04
This task follows the NLBSE'24 tool competition \cite{nlbse2024} and earlier
work on issue classification
\cite{kallis2020tickettagger,kallis2019tickettagger}.

The research questions are:

\begin{enumerate}
  \item Which text configuration performs best under training-only
  cross-validation?
  \item How well does the chosen model perform on the official test set,
  including differences across repositories and labels and comparison with
  the NLBSE'24 results?
  \item Which errors does the chosen model make most often, and how can the
  approach be improved?
\end{enumerate}

\section{Dataset and Invariants}

The NLBSE'24 tool competition provides 3,000 issues from five repositories. Each split contain % INTENTIONAL ESL 05
1,500 issues. The training set provide 100 examples for every repository and % INTENTIONAL ESL 06
label combination. The test set remain fixed for final evaluation. % INTENTIONAL ESL 07
It has the same balanced structure as the training set. No title is missing.
Two test issues have a missing description. Table~\ref{tab:dataset} summarizes
these checks.

\begin{table}[H]
\centering\small
\caption{Size and quality checks for the official dataset split.}
\label{tab:dataset}
\input{generated/dataset_summary.tex}
\end{table}

\section{Implementation Architecture}

The four notebooks are the main execution path. Each notebook call functions % INTENTIONAL ESL 08
from the project package; model logic is not duplicated inside the notebooks.
This sequence make the complete experiment easier to follow. % INTENTIONAL ESL 09

\begin{enumerate}
  \item Notebook 01 (\emph{Data Inspection}) calls the package data-loading,
  text-preparation, and inspection functions. It displays balance, missing
  values, duplicates, and exact overlaps.
  \item Notebook 02 (\emph{Chosen Model}) displays the executed five-fold
  comparison. Its optional rerun calls \texttt{run\_cross\_validation}, which
  fits all six candidates using training data only.
  \item Notebook 03 (\emph{Final Evaluation}) applies the selected
  configuration separately to five repositories. Its optional rerun calls the
  final-evaluation workflow.
  \item Notebook 04 (\emph{Results}) reads the final metrics and predictions to
  display repository scores, class scores, the confusion matrix, and common
  errors.
\end{enumerate}

For example, \texttt{build\_model\_text} prepares the issue text and
\texttt{run\_cross\_validation} performs the train-only model comparison.

The executed outputs are stored as JSON, CSV, and PNG files. For example,
\texttt{metrics.json} contains the final scores and \texttt{predictions.csv}
contains row-level true and predicted labels. The final-result directory also
contains confusion-matrix images. The notebooks read these files to display
tables and figures.

\section{Preprocessing, Data Inspection, and Protocol}

The preprocessing step join the issue title and description into one text field. % INTENTIONAL ESL 10
A missing description become an empty string. % INTENTIONAL ESL 11
Repeated whitespace is reduced to one space. These steps preserves code, % INTENTIONAL ESL 12
Markdown, URLs, identifiers, punctuation, and letter case. Repository name and
creation time are not model features.

The function first validates the five required columns. It then returns a new
Pandas series with the same row index as the input frame. It inserts
\texttt{[TITLE]} and \texttt{[BODY]} markers so the vectorizer can distinguish
the two fields. It does not lowercase, stem, remove stop words, or strip code.
Figure~\ref{fig:preprocess-code} shows the exact construction and the main
inspection operations.

\begin{figure}[H]
\centering
\includegraphics[width=0.98\linewidth]{generated/figures/code-preprocessing-inspection.png}
\caption{Text construction and per-split dataset inspection.}
\label{fig:preprocess-code}
\end{figure}

The inspection counts missing values and groups rows by repository and label.
It calculates text-length statistics and finds exact duplicate groups. For the
cross-split check, the code hashes the raw title-description pair with
SHA-256 after converting only a missing description to an empty string. It does
not normalize whitespace before this comparison. This keeps the overlap rule
strict and repeatable.

We also compare exact title-description pairs across the two splits. Three test
issues also appear in training. One of these pairs has different labels. The
main evaluation keeps all official test issues. A second evaluation removes
only these three overlaps to measure their effect.

\section{Chosen Model Selection}

The comparison include five text models and one dummy baseline. % INTENTIONAL ESL 13
The text models are two word TF--IDF logistic-regression configurations and
three word-and-character TF--IDF LinearSVC configurations. The dummy baseline
ignore the issue text and always predicts one class. % INTENTIONAL ESL 14

The word vectorizer uses one-word and two-word sequences. The character
vectorizer uses three-to-five-character sequences inside word boundaries. Both
use sublinear term frequency. Logistic regression is tested at $C=1$ and
$C=4$. LinearSVC is tested at $C=0.5$, $C=1$, and $C=2$.
Figure~\ref{fig:pipeline-code} shows how the word and character matrices are
joined before LinearSVC receives them.

\begin{figure}[H]
\centering
\includegraphics[width=0.90\linewidth]{generated/figures/code-candidate-pipeline.png}
\caption{Construction of the word-and-character TF--IDF LinearSVC pipeline.}
\label{fig:pipeline-code}
\end{figure}

Five-fold cross-validation divides each repository's training issues into five
parts. Each candidate use four parts for training and one part for validation. % INTENTIONAL ESL 15
This process repeats until every part has been used for validation. The five
folds gives one mean macro-F1 score for each repository and candidate. % INTENTIONAL ESL 16
We then average the five repository scores. The official test set is not read
during this comparison.

Each repository contains 300 training issues. One fold therefore fits on 240
rows and validates on 60 rows, with 20 validation rows for each label. The
vectorizers and classifier are created again inside every fold, so validation
text cannot influence the TF--IDF vocabulary. Six candidates, five
repositories, and five folds produce 150 fitted pipelines. Of these, 125 are
fits of the five text models and 25 are dummy fits. Figure~\ref{fig:fold-code}
shows the loop.

\begin{figure}[H]
\centering
\includegraphics[width=0.98\linewidth]{generated/figures/code-five-fold-comparison.png}
\caption{Training-only five-fold comparison for every repository and candidate.}
\label{fig:fold-code}
\end{figure}

Every candidate produces 25 macro-F1 values: five folds for each of five
repositories. The code first averages the five folds within each repository,
then averages the five repository means. It selects the highest global mean.
An exact tie is resolved by lower population standard deviation, lower
complexity rank, and candidate name, in that order.

Word-and-character TF--IDF with LinearSVC at $C=1$ achieved the highest mean
cross-validation macro-F1, 0.748. We therefore selected this configuration for
all five repositories. Table~\ref{tab:grid} reports every candidate, including
the dummy baseline.

\begin{table}[H]
\centering\small
\caption{Five-fold cross-validation results on the training set.}
\label{tab:grid}
\input{generated/candidate_grid.tex}
\end{table}

\section{Applying the Chosen Model}

The final stage read the selected candidate name from the saved % INTENTIONAL ESL 17
cross-validation result. It checks that this name matches a defined candidate
factory. It then loads only the
training split. For each repository, it create a fresh copy of the chosen % INTENTIONAL ESL 18
pipeline and fits it on all 300 training issues from that repository. This
produces five separate models with the same parameters but different TF--IDF
vocabularies and LinearSVC weights.

The five pipelines is saved as \texttt{.joblib} files. Their filenames, % INTENTIONAL ESL 19
repositories, and SHA-256 hashes are recorded in the training manifest. The
code verify that all five files exist and match their recorded hashes before % INTENTIONAL ESL 20
loading the test split. Figure~\ref{fig:final-application-code} shows the two
repository loops.

\begin{figure}[H]
\centering
\includegraphics[width=0.98\linewidth]{generated/figures/code-final-model-application.png}
\caption{Training five repository models and applying each model to its test rows.}
\label{fig:final-application-code}
\end{figure}

The official test split is loaded once after training and verification. Each
repository model receives the 300 test issues from the same repository. The
primary policy keeps all 1,500 predictions. The leakage-sensitive policy keeps
1,497 predictions after removing the three exact overlaps. Each stored row
contains the repository, true label, predicted label, text hash, and overlap
flag. No test label is passed to \texttt{fit()} or used to change the pipeline.

\section{Final Chosen-Model Evaluation}

Our chosen model is trained once on all training issues from each repository.
It achieves a cross-repository macro-F1 of 0.754385 on the official test set.
The score is the arithmetic mean of the five repository macro-F1 scores.
Tables~\ref{tab:repositories} and~\ref{tab:classes} show where performance
varies by repository and label.

\begin{table}[H]
\centering\small
\caption{Repository-level macro metrics for our chosen model.}
\label{tab:repositories}
\input{generated/chosen_repositories.tex}
\end{table}

\begin{table}[H]
\centering\small
\caption{Global per-class metrics for both evaluation policies.}
\label{tab:classes}
\input{generated/chosen_classes.tex}
\end{table}

\begin{table}[H]
\centering\small
\caption{Primary and leakage-sensitive policy summary.}
\label{tab:policies}
\input{generated/policy_summary.tex}
\end{table}

\begin{figure}[H]
\centering
\includegraphics[width=0.78\linewidth]{generated/figures/cross-validation-candidates.png}
\caption{Zoomed comparison of the five text models. Error bars show
population standard deviations; the dummy remains in Table~\ref{tab:grid}.}
\label{fig:cv}
\end{figure}

\section{Comparison with NLBSE'24 Results}

The NLBSE'24 tool competition reports SetFit scores of 0.827 (headline) and
0.824 (detailed). Their RoBERTa and fastText scores are 0.792 and 0.718. Our
chosen model achieved 0.754 on the official test set.

\begin{table}[H]
\centering\small
\caption{Our results compared with the NLBSE'24 results.}
\label{tab:baselines}
\input{generated/baseline_comparison.tex}
\end{table}

\section{Error Analysis}

The largest error occur when questions are predicted as features. % INTENTIONAL ESL 21
The next largest error is predicting bugs as questions. Table~\ref{tab:errors}
shows the ten most frequent repository-specific errors. Figure~\ref{fig:confusion}
shows all predictions across the five repositories.

\begin{table}[H]
\centering\small
\caption{Most frequent primary-policy confusion pairs.}
\label{tab:errors}
\input{generated/error_analysis.tex}
\end{table}

\begin{figure}[H]
\centering
\includegraphics[width=0.62\linewidth]{generated/figures/chosen-model-confusion-matrix.png}
\caption{Aggregate primary-policy confusion matrix generated from final metrics.}
\label{fig:confusion}
\end{figure}

\section{Analysis and Improvements}

The error table shows two main problems. The model confuses questions with
features because both labels contain requests for a change or a new capability.
It also confuses bugs with questions when an issue describes a failure and asks
for an explanation at the same time. Word and character features capture terms
such as ``error'', ``support'', and ``request'', but they do not reliably capture
the purpose of the sentence.

Performance also changes by repository. Bitcoin is the weakest repository with
macro-F1 of 0.701. Its issue language is more technical and its reports contain
shorter context. React is stronger at 0.807, where issue descriptions provide
more repeated vocabulary. This suggests that repository conventions affect the
task even when every repository has the same label counts.

These error patterns shows where the current features are insufficient. % INTENTIONAL ESL 22

The chosen model is one global configuration. It is simple and fast, but it
cannot adapt its vocabulary or decision boundary to each repository. The main
improvements are repository-specific tuning, more careful treatment of short
titles, and error-driven review of question--feature and bug--question pairs.
Character features can be expanded for identifiers and code fragments. A
sentence-embedding model can be tested later if a stronger semantic comparison
is required, but it is outside the chosen-model result reported here.
The proposed improvements focuses on the observed confusion pairs. % INTENTIONAL ESL 23

The three exact train/test overlaps have little effect. Removing them changes
macro-F1 from 0.754385 to 0.754642, a difference of 0.000257. The reported
result therefore does not depend materially on these rows.

\section{Run Instructions and Conclusion}

The four numbered notebooks are the recommended execution path. For local use,
create a Python 3.11 environment and call its interpreter directly. PowerShell
environment activation is not required. The same interpreter installs the
dependencies and starts JupyterLab:

\begin{verbatim}
py -3.11 -m venv .venv
.\.venv\Scripts\python.exe -m pip install -r requirements.txt
.\.venv\Scripts\python.exe -m jupyter lab
\end{verbatim}

After JupyterLab opens, run Notebooks 01, 02, 03, and 04 in numerical order.
By default, model training is off. Notebook 02 sets
\texttt{RERUN\_CROSS\_VALIDATION=False}, and Notebook 03 sets
\texttt{RERUN\_FINAL\_EVALUATION=False}. With these defaults, the notebooks
load and display the saved results. Set the relevant switch to \texttt{True}
only when the corresponding training stage must be repeated.

The same notebook workflow also runs in Colab. Upload the complete project
folder to Google Drive, open the notebooks from \texttt{notebooks/} in
numerical order, and select \emph{Runtime $>$ Run all}. The README lists
optional command-line checks, but they are not required for the notebook
presentation workflow.

Our chosen model achieved 0.754. The NLBSE'24 tool competition reports 0.827
for SetFit. The three exact train/test overlaps changed the chosen-model result
by only 0.000257. The experiment selected one configuration, measured its test
performance, and identified its main errors and possible improvements. These
results answer the three research questions without changing the official data
split.

\clearpage
\bibliographystyle{plain}
\bibliography{references}

\end{document}
